% Lernkarten -- Analysis (Modul 61211), Lektion 6
% Quelle: eigene handschriftliche Notizen (Fernuni Notizbuch > 61211 Analysis, Kapitel 6)
% Format: \lernkarte{Vorderseite (roter Titel)}{Rueckseite (weisser Inhalt)}

% --- Notiz-Seite 6.1 "Kurvenbegriff" -----------------------------------
\lernabschnitt{6.1 Kurvenbegriff}

\lernkarte[61211-6-01]{Parametrisierte Kurve}{%
  $I$ ein Intervall aus $\mathbb{R}$,
  \[ \varphi : I \longrightarrow \mathbb{R}^n \]
  ist eine parametrische Kurve. \\
  Spur von $\varphi := \mathrm{Bild}(\varphi)$.
}

\lernkarte[61211-6-02]{Äquivalenzrelation (Umparametrisierung)}{%
  Streng monoton wachsende Funktion $f$ mit $f(I) = J$ und
  $\varphi = \psi \circ f$, \; $\varphi : I \to \mathbb{R}^n$, \; $\psi : J \to \mathbb{R}^n$.
  \begin{enumerate}[label=(\roman*),topsep=2pt,itemsep=1pt,leftmargin=*]
    \item $\varphi \sim \varphi$
    \item $\varphi \sim \psi \Rightarrow \psi \sim \varphi$
    \item $\varphi \sim \psi,\ \psi \sim \chi \Rightarrow \varphi \sim \chi$
  \end{enumerate}
}

\lernkarte[61211-6-03]{Eigenschaften äquivalenter parametrisierter Kurven}{%
  \begin{enumerate}[topsep=2pt,itemsep=1pt,leftmargin=*]
    \item[(0)] $[\varphi] = [\psi]$
    \item[(i)] $\mathrm{Spur}(\varphi) = \mathrm{Spur}(\psi)$
    \item[(ii)] $t_1 < t_2 \Rightarrow \exists\, \tau_1 < \tau_2$ mit
          $\varphi(t_1) = \psi(\tau_1)$ und $\varphi(t_2) = \psi(\tau_2)$
    \item[(iii)] $\varphi(\text{Endpunkt}) = \psi(\text{Endpunkt})$
  \end{enumerate}
}

\lernkarte[61211-6-04]{Kurve}{%
  Eine Kurve $W$ ist eine Äquivalenzklasse $[\varphi]$ parametrisierter Kurven.
}

% --- Notiz-Seite 6.3 "Kurvenintegrale" ---------------------------------
\lernabschnitt{6.3 Kurvenintegrale}

\lernkarte[61211-6-05]{Kurvenintegral für skalare Funktion}{%
  Voraussetzung: $W$ ist glatt, $\varphi$ auch glatt, $f : |W| \to \mathbb{R}$
  stetig. So heißt
  \[ \int_W f\, ds := \int_a^b f(\varphi(t))\, \|\dot{\varphi}(t)\|_2\, dt \]
  das Kurvenintegral.
}

\lernkarte[61211-6-06]{Kurvenintegral für Vektorfunktionen}{%
  Voraussetzung: glatte Kurve, glatte Parameterdarstellung.
  \[ \int_W f \cdot dx := \int_a^b f(\varphi(t)) \cdot \dot{\varphi}(t)\, dt \]
  oder
  \[ \int_W (f_1\, dx_1 + \ldots + f_n\, dx_n) \]
}

% --- Notiz-Seite 6.4 "Stammfunktion" -----------------------------------
\lernabschnitt{6.4 Stammfunktion}

\lernkarte[61211-6-07]{Gebiet}{%
  Nicht leere, zusammenhängende, offene Teilmenge von $\mathbb{R}^n$.
}

\lernkarte[61211-6-08]{Stammfunktion}{%
  $G$ ist ein Gebiet, $f : G \to \mathbb{R}^n$, $F : G \to \mathbb{R}$. \\
  Wenn ${}^t F'(x) = f(x)$, dann heißt $F$ eine Stammfunktion von $f$.
}

\lernkarte[61211-6-09]{Potentialfeld}{%
  $U := -F$ heißt Potenzialfeld zum Vektorfeld $f$.
}

\lernkarte[61211-6-10]{Gradientenfeld}{%
  $f$ ist das Gradientenfeld, erzeugt vom Potenzial $U := -F$.
}

\lernkarte[61211-6-11]{Integritätsbedingungen}{%
  Besitzt $f$ eine Stammfunktion, so gelten
  \[ D_i f_j = D_j f_i \]
}

\lernkarte[61211-6-12]{Wegunabhängigkeit}{%
  Voraussetzung: $W$ stückweise glatte Kurve, mit Anfangs- und Endpunkt
  $a, b$ und $|W| \subseteq G$. Dann
  \[ \int_W f \cdot dx = F(b) - F(a) \]
  Also spielt der Verlauf keine Rolle.
}

\lernkarte[61211-6-13]{Satz Stammfunktion (Existenz)}{%
  Sei $G$ ein Gebiet, $f$ stetig differenzierbar, erfüllt die
  Integritätsbedingung. Ist $G$ sternförmig, so besitzt $f$ eine
  Stammfunktion $F$.
}

\lernkarte[61211-6-14]{Sternförmig}{%
  \[ \exists\, a \in G \ \forall\, u \in G : \ S(a, u) \subseteq G \]
}

% --- Notiz-Seite 6.5 "Flächen und Volumenberechnung" -------------------
\lernabschnitt{6.5 Flächen und Volumenberechnung}

\lernkarte[61211-6-15]{Flächeninhalt}{%
  \[ \text{Fläche} = F := \left\{ \tbinom{x}{y} \in \mathbb{R}^2 \ \middle|\
     x \in I \text{ und } 0 \le y \le f(x) \right\} \]
  \[ \text{Flächeninhalt} = \lambda_2(F) := \int_I f \]
}

\lernkarte[61211-6-16]{Körper zwischen zwei Graphen}{%
  \begin{align*}
    \lambda_3(K) &= \int_a^b \lambda_2(F(x))\, dx \\
                 &= \int_a^b \left( \int_{\alpha(x)}^{\beta(x)}
                    (f(x,t) - g(x,t))\, dt \right) dx
  \end{align*}
}
